% =============================================================================
% preamble.tex — Преамбула диссертации (XeLaTeX)
% =============================================================================

% --- Шрифты и язык (XeLaTeX) ---
\usepackage{fontspec}
\setmainfont{Times New Roman}
\setsansfont{Arial}
\setmonofont{Courier New}

\usepackage{polyglossia}
\setdefaultlanguage{russian}
\setotherlanguage{english}

% --- Геометрия страницы (ГОСТ) ---
\usepackage[
  a4paper,
  left=30mm,
  right=15mm,
  top=20mm,
  bottom=20mm
]{geometry}

% --- Межстрочный интервал 1.5 ---
\usepackage{setspace}
\onehalfspacing

% --- Абзацный отступ 1.25 см ---
\setlength{\parindent}{1.25cm}

% --- Математика ---
\usepackage{amsmath,amssymb}

% --- Графика и рисунки ---
\usepackage{graphicx}
\graphicspath{{images/}}
\usepackage{float}

% --- Таблицы ---
\usepackage{longtable}
\usepackage{booktabs}
\usepackage{array}
\usepackage{multirow}
\usepackage{tabularx}

% --- Списки ---
\usepackage{enumitem}
\setlist{nosep, leftmargin=\parindent}

% --- Подписи к рисункам и таблицам ---
\usepackage{caption}
\captionsetup{
  labelsep=endash,
  justification=centering,
  font={onehalfspacing},
  skip=6pt
}
\captionsetup[figure]{name=Рисунок}
\captionsetup[table]{name=Таблица}

% --- Оглавление ---
\usepackage{titletoc}
\usepackage{tocloft}
\renewcommand{\cfttoctitlefont}{\hfill\large\bfseries}
\renewcommand{\cftaftertoctitle}{\hfill}
\renewcommand{\cftdotsep}{1}
\renewcommand{\cftsecleader}{\cftdotfill{\cftdotsep}}
\setlength{\cftbeforesecskip}{3pt}

% --- Заголовки глав и разделов ---
\usepackage{titlesec}

% Глава: "ГЛАВА 1. НАЗВАНИЕ" — по центру, жирный, прописными
\titleformat{\section}[block]
  {\normalfont\large\bfseries\centering}
  {\thesection.}{0.5em}{\MakeUppercase}
\titlespacing*{\section}{0pt}{24pt}{12pt}

% Раздел: "1.1. Название" — с абзацного отступа, жирный
\titleformat{\subsection}[block]
  {\normalfont\normalsize\bfseries}
  {\thesubsection.}{0.5em}{}
\titlespacing*{\subsection}{\parindent}{18pt}{6pt}

% Подраздел: "1.1.1. Название" — с абзацного отступа, жирный
\titleformat{\subsubsection}[block]
  {\normalfont\normalsize\bfseries}
  {\thesubsubsection.}{0.5em}{}
\titlespacing*{\subsubsection}{\parindent}{12pt}{6pt}

% --- Нумерация страниц (внизу по центру) ---
\usepackage{fancyhdr}
\pagestyle{fancy}
\fancyhf{}
\fancyfoot[C]{\thepage}
\renewcommand{\headrulewidth}{0pt}
\renewcommand{\footrulewidth}{0pt}

% plain-стиль тоже с номером внизу по центру
\fancypagestyle{plain}{
  \fancyhf{}
  \fancyfoot[C]{\thepage}
  \renewcommand{\headrulewidth}{0pt}
}

% --- Листинги кода ---
\usepackage{listings}
\lstset{
  basicstyle=\small\ttfamily,
  breaklines=true,
  frame=single,
  numbers=left,
  numberstyle=\tiny,
  tabsize=2,
  captionpos=b,
  extendedchars=true,
  inputencoding=utf8,
  showstringspaces=false
}

% --- Библиография (ГОСТ) ---
\usepackage[
  backend=biber,
  style=gost-numeric,
  sorting=none,
  language=auto,
  autolang=other
]{biblatex}
\addbibresource{references.bib}

% --- Гиперссылки ---
\usepackage{hyperref}
\hypersetup{
  colorlinks=true,
  linkcolor=black,
  citecolor=black,
  urlcolor=blue,
  pdfencoding=auto
}

% --- Приложения ---
\usepackage[title]{appendix}

% --- Для удобства: команда \todo ---
\newcommand{\todo}[1]{\textcolor{red}{\textbf{TODO:} #1}}

% --- Счётчики для сквозной нумерации рисунков и таблиц ---
\usepackage{chngcntr}
